%% kapa: exam=1 round=26 year=2015 period=2 form=A subject=accounting subject_ko=회계학 seq=7 number=47
%% source: Q-Net 감정평가사 기출문제 (cst00309) artlSeq=5138001 file=extracted/1008/5138001/제26회_감정평가사 1차_기출/제26회_감정평가사 1차_기출/2교시_감평_A형.pdf
%% conversion: kordoc 4.24.51 -> tools/gen1.py (auto); review=proofread-sonnet; flags=none
\begin{question}{47}
20×1년 1월 1일에 설립된 (주)감평은 20×1년 말에 확정급여제도를 도입하였다. 확정급여채무 계산 시 적용한 할인율은 연 10\%이며, 20×1년 이후 할인율의 변동은 없다. 다음 자료를 이용하여 계산된 20×2년 순확정급여부채는?
\begin{framedtext}
○ 20×1년 말 확정급여채무의 장부금액은 ￦30,000이다. \par
○ 20×1년 말 사외적립자산에 ￦20,000을 현금으로 출연하였다. \par
○ 20×2년 말 퇴직한 종업원에게 ￦1,000의 현금이 사외적립자산에서 지급되었다. 당기 근무원가는 ￦40,000이다. \par
○ 20×2년 말 사외적립자산에 ￦30,000을 현금으로 출연하였다. \par
○ 20×2년 말 현재 사외적립자산의 공정가치는 ￦65,000이다. \par
○ 20×2년 말 보험수리적 가정의 변동을 반영한 확정급여채무의 현재가치는 ￦80,000이다.
\end{framedtext}
\begin{choices}
  \item ￦15,000
  \item ￦20,000
  \item ￦25,000
  \item ￦65,000
  \item ￦80,000
\end{choices}
\end{question}
\answer{1}
